
\chapter{Experiments}\label{ch:experiments}

This chapter reports the experimental evaluation of the proposed runtime. The goal is twofold: to confirm that a single \texttt{--security-scan} pass produces working seccomp, AppArmor, and capability profiles (RQ1), and to measure whether scan-then-enforce prevents command-injection exploitation and improves containment relative to stock \texttt{runc} (RQ2). The evaluation has an empirical part, run on the \texttt{runc-hardened-test} security laboratory, and an analytical part that positions the runtime's performance against sandboxed alternatives using published measurements. Section~\ref{sec:exp-setup} describes the experiment and its environment; Section~\ref{sec:exp-results} reports the results.

\section{Experimental Setup}\label{sec:exp-setup}

\textbf{Experiment and implemented features.} For the proof of concept, a small intentionally vulnerable web application was developed. The application contains improper input validation that allows command injection through an exposed endpoint. Around it, scripts were written to reproduce the full baseline\,/\,scan\,/\,enforce\,/\,attack cycle. The laboratory implements three feature groups, each tied to a question the experiment must answer:

\begin{itemize}
    \item \textbf{Web application functionality} --- the most likely scenario for a containerised workload, used as the legitimate traffic that drives the scan;
    \item \textbf{System monitoring, calendar rendering, and upload functionality} --- exercised to check that the generated profiles do not break the application after enforcement;
    \item \textbf{Vulnerable command-execution endpoints} --- used to check whether the enforced profiles block exploitation of the injected command.
\end{itemize}

\textbf{Environment.} The following specifications were used:

\begin{itemize}
    \item Ubuntu 24.04 LTS;
    \item cgroup v2 enabled;
    \item AppArmor enabled;
    \item root privileges for \texttt{runc}, AppArmor profile loading, and BPF tracing;
    \item stock \texttt{runc} v1.0.0-rc95 as the baseline runtime.
\end{itemize}

\textbf{Reproducibility.} The complete laboratory setup, the vulnerable application, and step-by-step instructions for replaying the experiment are published in the \texttt{runc-hardened-test} repository: \url{https://github.com/dpttk/runc-hardened-test} (see Appendix~\ref{app:repos}).

\section{Experimental Results}\label{sec:exp-results}

The experiment was successful: the proposed runtime generated profiles that prevented command injection on the vulnerable endpoint while preserving the application's legitimate functionality.

\subsection{Profile generation (RQ1)}

A single \texttt{runc run --security-scan} pass, followed by scripted legitimate traffic, produced the expected artefacts under \texttt{oci-bundle/generated/}: \texttt{seccomp.json} with \texttt{defaultAction=SCMP\_ACT\_ERRNO}, an AppArmor profile learned in complain mode, and \texttt{capable-bpfcc.log} from the cgroup-wide BPF trace. Finalisation embedded the seccomp allow-list and narrowed \texttt{process.capabilities} in \texttt{config.json} without widening privileges relative to the trace. Table~\ref{tab:profquality-lab} summarises the artefact families.

\begin{table}[ht]
\centering
\caption[Profile artefacts produced by scan.]{Profile artefacts produced by \texttt{--security-scan} on the runtime-security lab workload.}
\label{tab:profquality-lab}
\begin{tabularx}{\textwidth}{>{\raggedright\arraybackslash}p{4.2cm} X}
\toprule
\textbf{Artefact} & \textbf{Role} \\
\midrule
\texttt{generated/}\\\texttt{seccomp.json} & Deny-by-default syscall allow-list (\texttt{oci-seccomp-bpf-hook}) \\
\texttt{generated/}\\\texttt{apparmor.profile} & MAC rules from complain-mode audit \\
\texttt{generated/}\\\texttt{capable-bpfcc.log} & Observed \texttt{cap\_capable} events (cgroup-wide) \\
\texttt{config.json} (post-finalise) & Narrowed \texttt{process.capabilities}, embedded seccomp \\
\bottomrule
\end{tabularx}
\end{table}

\subsection{Containment effectiveness (RQ2)}

After scan-then-enforce, the same injection payloads from the baseline run were replayed against the hardened configuration. Under stock \texttt{runc} they succeeded; under the hardened runtime the injected commands did not execute.

\textbf{Seccomp prevented command injection.} The decisive control was the generated seccomp allow-list. Because the scan never observed process-creation syscalls on the attack path, the enforced profile denied \texttt{execve}\,/\,\texttt{execveat}: the injected \texttt{sh -c} chain could not spawn \texttt{bash} or fork any helper, so the command-injection primitive failed at the syscall boundary. The same denial also stopped reconnaissance, sensitive-file reads, the reverse shell, and the mount probe. This is the layer that \texttt{runc} itself enforces.

\textbf{AppArmor preserved legitimate file access.} In parallel, the learned AppArmor profile allowed exactly the files the application needs: calendar storage under \path{/var/lib/sysmon-cal} and the application log paths remained readable and writable, so calendar rendering, uploads, and logging kept working after enforcement.

\textbf{Defence in depth beyond the exec block.} As a controlled ablation, \texttt{execve} was manually re-enabled in the seccomp profile. That restored a shell (RCE), but most follow-up actions still stayed blocked: the rest of the syscall allow-list, the narrowed capability set, and the AppArmor rules denied access to the commands and host files the attacker tried to use. Under normal enforcement, injection is stopped at \texttt{execve}; the ablation shows that the remaining layers still limit damage if \texttt{execve} were ever allowed.

The following payloads succeeded under stock \texttt{runc} but were blocked after scan-then-enforce (including the injection itself):

\begin{itemize}
    \item \textbf{Reconnaissance} --- \texttt{id; hostname; uname -a};
    \item \textbf{Sensitive read} --- \texttt{cat /etc/shadow};
    \item \textbf{Reverse shell} --- \texttt{bash -i >\& /dev/tcp/\ldots};
    \item \textbf{Mount probe} --- \texttt{mount -t proc proc /mnt};
    \item \textbf{Persistence} --- \texttt{chmod u+s /bin/bash}.
\end{itemize}

Legitimate ping (\texttt{/diag/ping?host=127.0.0.1}) and calendar upload (\texttt{POST /calendar/upload}) remained functional.

\subsection{Performance: analytical comparison}\label{sec:exp-analytical}

This thesis does not run its own performance benchmarks; the comparison below is analytical and draws on published measurements. The relevant question is the steady-state overhead each confinement approach imposes relative to stock \texttt{runc}, which is the natural baseline because all three options either are \texttt{runc} or replace it.

\textbf{Stock \texttt{runc}.} As a thin shared-kernel runtime, \texttt{runc} runs the workload directly on the host kernel and is treated as the near-native baseline (0\% reference overhead); its weakness is isolation, not speed~\cite{wang2022,volpert2024}.

\textbf{gVisor.} gVisor interposes a user-space kernel (Sentry) that intercepts and re-implements guest syscalls outside the host kernel. This buys strong isolation but adds a per-syscall user-space round trip, so syscall- and I/O-heavy workloads pay the most: reported overheads reach roughly 50--100\% relative to default \texttt{runc}~\cite{viktorsson2020,wang2022}.

\textbf{Kata Containers.} Kata runs each container inside a lightweight virtual machine with its own kernel. The isolation boundary is hardware-assisted, but the microVM adds memory cost and a startup penalty, with steady-state overheads commonly reported around 20--40\%~\cite{viktorsson2020,volpert2024}.

\textbf{This work (hardened \texttt{runc}).} The proposed runtime keeps the shared-kernel model and only adds enforcement of the generated profiles. At steady state this is the cost of ordinary seccomp filtering --- a per-syscall BPF check --- plus AppArmor path checks and a smaller capability set, none of which changes the execution model. Until dedicated benchmarks are run (left to future work, Chapter~\ref{ch:conclusion}), the working assumption is therefore that the runtime's overhead relative to stock \texttt{runc} stays at the level of normal seccomp operation, i.e.\ far below the sandbox figures above. The one-time scan pass carries a higher eBPF-tracing cost, but that is a build-time activity, not a steady-state property of the enforced container. Table~\ref{tab:perf-ch5} summarises the positioning.

\begin{table}[ht]
\centering
\caption{Analytical performance positioning relative to stock \texttt{runc}. Sandbox figures are taken from published measurements; the hardened row is an assumption pending dedicated benchmarks.}
\label{tab:perf-ch5}
\resizebox{\textwidth}{!}{%
\begin{tabular}{@{}p{3.4cm}p{4.6cm}p{4.2cm}@{}}
\toprule
\textbf{Runtime} & \textbf{Isolation model} & \textbf{Steady-state overhead vs.\ \texttt{runc}} \\
\midrule
Stock \texttt{runc} & Shared host kernel & Baseline (near-native) \\
This work (hardened \texttt{runc}) & Shared kernel + generated seccomp/AppArmor/capabilities & $\approx$ seccomp filtering cost (assumed; not yet measured) \\
gVisor & User-space kernel (Sentry) & $\sim$50--100\%~\cite{viktorsson2020,wang2022} \\
Kata Containers & Per-container microVM & $\sim$20--40\%~\cite{viktorsson2020,volpert2024} \\
\bottomrule
\end{tabular}%
}
\end{table}
